\section{Preliminary Mathematical Properties}
Proofs to theorems mentioned are provided in the Appendix \autoref{subsec:app-prelim_math}.

	\subsection*{Programmatic Notation}
	We follow these conventions of nomenclature, in both the report and the code as far as possible, to quickly understand the domain of variables, and in some ways also their interpretation.
	\begin{enumerate}
		\item $x \in [a, a+d]$
		\item $z :\equiv\frac{x-a}{d} \in [0,1]$
		\item $u :\equiv z^s\in [0,1]$ also.
		\item $\theta :\equiv \pi u \equiv \pi z^s \in [0,\pi]$
		\item $\omega := \theta - \frac{\pi}{2} \equiv \pi z^s - \frac{\pi}{2} \in [-\frac{\pi}{2}, \frac{\pi}{2}]$
		\item $r = \pm 1$
	\end{enumerate}


\subsection{Physical Interpretation}
Similar to \cite{gilbert1893moon} -- which imagines intensity (interpreted as probability density) of crater activity on a circular planet as proportional to Gilbert's Sine density $\sin\theta$ -- the Standard \sindist{} $\Sinu(s,k)$ can also be physically interpreted as the intensity of crater activity on a (possibly non-circular) planetary surface given by the radial curve:
\begin{equation}
	r(z) = r_0 \exp\left(\int_0^{\pi z^s} \tan(t - \sin^{-1}\sin^k t) \d t\right)
\end{equation}
where the curve $r(z)$ is parametrized by $z \in [0,1]$ such that $\theta=\pi z^s$ is the angle the location $r(z)$ subtends at the origin, and $r_0$ is the initial radius at $z=0$. Details are furnished in \autoref{rmk:planet}.


	\subsection{Transformation}
	\begin{thm}[Expression as a morphed arcsin Beta distribution] \label{thm:transform}
		$Y \sim \Beta(\frac{k+1}{2}, \frac{1}{2})$ and $R \sim \Ber (\frac 1 2)$, then
		\[Z = \left[ R \sin^{-1} \sqrt Y + (1-R) (\pi - \sin^{-1} \sqrt Y) \right]^\frac{1}{s}\]
		follows a $\Sinu(s,k)$ distribution.
	\end{thm}



\begin{thm}[Joint density expression]\label{thm:jtdens}
	For $X_1\dots X_n \iidsim \Sinu(a,d,s,k)$, their joint density can be expressed as:
	\begin{equation}
		f(x_1\dots x_n) = \frac{1}{A^n(s,k) 2^{nk} d^n } \left( \sum_{r_1=\pm1}\cdots \sum_{r_n = \pm 1} \left[\prod_{j=1}^n r_j\right] \cdot \exp\left[\sum_{j=0}^n \i \pi r_j \left(\frac{x_j -a}{d}\right)^s\right]  \right)^k
	\end{equation}
\end{thm}

\subsection{Concentration Inequalities}
\begin{thm}[\sindist{} is sub-Gaussian]
	$X \sim \Sinu (a,d,s,k)$ is sub-Gaussian with parameter $\frac{d^2}{4}$.
\end{thm}
\begin{proof}
	$\because a \le X \le a+d \quad\text{a.s.}$, $\therefore \E \left[ \e^{\lambda (X - \E X)} \right] \le \exp \left( -\frac{\lambda d^2}{8} \right)$ by Hoeffding's lemma.
	
	Note that because it is sub-Gaussian, it is also sub-Exponential.
\end{proof}


\subsection{Normalizing constant for the symmetric case}
\begin{thm}[Normalizing constant, symmetric case]\label{thm:norm_sym}
	\begin{align}
		A(1,k) \equiv \int_0^1 \sin^k(\pi t) \d t &= \frac{2^k}{\pi} \frac{\Gamma^2\left( \frac{k+1}{2} \right)}{\Gamma\left( k+1\right)} \label{eq:area1}\\
		&= \frac{1}{\sqrt \pi} \frac{\Gamma\left(\frac{k+1}{2}\right)}{\Gamma\left(\frac{k}{2} + 1\right)} \label{eq:area2}
	\end{align}
\end{thm}


\subsection{Incomplete integral of the kernel for the Symmetric case}
\begin{thm}[Incomplete integral of the kernel for the symmetric case] \label{thm:cdf_sym}
	When $k\in \ZZ^+$, $s=1$, we have --
	\begin{align}
		\int_0^z \sin^k\pi t \, dt &=
		\frac{1}{2^{2n}} {2n \choose n} z + \frac{(-1)^n}{\pi 2^{2n-1}}\sum_{k=0}^{n-1} (-1)^k {2n \choose k} \frac{\sin 2(n-k)\pi z}{2(n-k)} && [\text{$k$ even}]\\
		&= \frac{1}{\pi 2^{2n}}(-1)^n \sum_{k=0}^n (-1)^k {2n+1 \choose k} \frac{1 - \cos (2n+1-2k)\pi z}{2n+1-2k}  &&  [\text{$k$ odd}]
	\end{align}
\end{thm}


\subsection{Limiting to Normality}
\begin{thm}\label{thm:normality}
	For the symmetric about 0 \sindist{} $\Sinu\left(-\frac{d}{2}, d, 1, k\right)$, i.e. with $a = -\frac{d}{2}$ and $s=1$, setting $d \equiv d(k)= \pi \sqrt k$ and letting $k \uparrow \infty$, we have:
	\[ \psi_k(x) \to \phi(x)\]
	where $\phi(x)$ is the PDF of a Standard Normal distribution.
\end{thm}

\subsection{Density of Sum}
\begin{thm}
	Density of Sum
\end{thm}
\begin{proof}
	Assume
	\begin{align*}
		content...
	\end{align*}
\end{proof}


